% The simplicial cone σ = Cone(e1, e2, e1 + e2 + 3e3) in R^3, in a projection
% where e1 points down-left, e2 right, and e1 + e2 + 3e3 up out of the
% e1-e2 plane. The ray through e2 lies behind the cone and is dotted.
\documentclass[tikz,border=2pt]{standalone}
\usepackage{dzg-tikz}

\begin{document}
\begin{tikzpicture}[scale=1.3]
  \coordinate (O)  at ( 0,    0  );
  \coordinate (E1) at (-2.4, -1.6);
  \coordinate (E2) at ( 4.0,  0  );
  \coordinate (E3) at ( 1.6,  5.6);

  % Facets, back to front.
  \fill[dzg grid] (O) -- (E1) -- (E3) -- cycle;
  \fill[dzg blob] (O) -- (E1) -- (E2) -- cycle;
  \fill[dzg grid!60] (O) -- (E2) -- (E3) -- cycle;

  % Rays, continuing past the generators to show the cone is unbounded.
  \draw[fan ray, dotted] (O) -- ($(O)!1.15!(E2)$);
  \draw[fan ray] (O) -- ($(O)!1.15!(E1)$);
  \draw[fan ray] (O) -- ($(O)!1.15!(E3)$);

  \node[distinguished point, label={below left:$0$}] at (O) {};
  \node[distinguished point, label={below left:$e_1$}] at (E1) {};
  \node[distinguished point, label={below right:$e_2$}] at (E2) {};
  \node[distinguished point, label={above left:$e_1 + e_2 + 3e_3$}] at (E3) {};
\end{tikzpicture}
\end{document}
